%!TEX root = ../Polito PhD thesis template.tex
% ******************************* Chapter 2 ************************************
% PART I OPENER ("cappello"): theory and common ground for the LFR works.
% Only theory and general framework, as agreed for the Part II opener: the
% design-specific material, the generation strategies, the SPH algorithm and
% the optimization of the energy mesh belong to Chapters 3 and 4.
%
% Notation: Part I uses \vec{r} and \vec{\Omega}, as Chapter 3 does; Part II
% uses \mathbf{r}. Harmonising the two is a cosmetic decision left open.
% The angular flux is \phi and the total flux, integrated over the angle,
% is \Phi; \psi is left free for the weighting function. Chapter 3 still
% writes \phi for the total flux, which is the notation of the published
% paper.
%
% Sources: standard transport theory (Bell and Glasstone); the definitions of
% the group constants follow the introductions of the two works of this Part,
% materials/enanching_multigroup_LFR/SPH_paper_NS_E/main_rev_v1_clean.tex and
% materials/enanching_multigroup_LFR/lfr_energy_grid/paper/paper.tex.

\chapter{Multi-Group Data for Lead-Cooled Fast Reactor Analysis}\label{ch:part1-opener}

\section{Motivation and scope}

The design of a reactor and the analysis of its safety require the neutron field
to be computed many times: for the configurations of the control system that the
plant will pass through, for the states of the fuel that follow one another
along the cycle, and, in the transient analyses for which those calculations are
ultimately intended, for many instants in time. A Monte Carlo (MC) calculation
answers each of those questions with the fewest approximations, since it treats
energy and angle as continuous variables and represents the geometry of the
system in full detail. The general-purpose codes that implement it, among which
MCNP~\cite{mcnp63}, Serpent~2~\cite{Serpent_code_description} and
OpenMC~\cite{openmc} are the most widely used in reactor physics, are therefore
the reference against which every approximate model is measured. Their cost per
configuration is what confines them to a limited number of reference states.

Deterministic solvers return the same quantities in a time smaller by orders of
magnitude, and it is with them that parametric studies and transient analyses
are usually carried out. The separation is one of cost rather than of principle,
since time-dependent Monte Carlo modes have been developed as well, in Serpent~2
among other codes~\cite{serpentDynamic}; but a study that must repeat the same
calculation over hundreds of configurations still turns to a deterministic
solver.

The price of that speed is paid in the form in which the nuclear data have to be
supplied. A deterministic solver does not read the continuous-energy cross
sections: it reads a condensed and homogenized version of them, the multi-group
constants, that is a set of cross sections defined over a finite number of
energy intervals and over regions of the core within which the material is
treated as uniform. Producing those constants is a calculation in its own right,
and the accuracy of everything computed afterwards is bounded by their quality. % CV: mi sto rendendo conto che lo stile è molto lavativo, assicurati che non sia cosi su tutto il resto della tesi. Voglio che lo stile di scrittura sia ispirato a quello che trovi dentro la nuova cartella che ti ho generato: phd_tesis_sampel_nicolo (tesi del mio supervisor).

This Part examines two of the degrees of freedom that decide that quality, for
Lead-cooled Fast Reactors (LFRs). Chapter~\ref{ch:sph-lfr} concerns the
equivalence between the heterogeneous system and its homogenized representation,
restored by correcting the constants themselves. Alongside the correction, it
compares the routes by which the constants can be generated, a single full-core
Monte Carlo calculation against a set of mini-core models representative of the
different regions of the system, and weighs the accuracy each route delivers
against the computational time it costs when several configurations of the core
have to be analyzed.
Chapter~\ref{ch:energy-grid}
concerns the energy mesh on which the condensation is performed, whose
boundaries are treated as the unknown of an optimization. The two works share
the same scheme, in which the Monte Carlo code Serpent~2
\cite{Serpent_code_description,Serpent_homog_technique} generates both the
multi-group constants and the reference solution against which the deterministic
result is judged. Since the same continuous-energy nuclear data serve both
sides, the comparison measures the multi-group model and not the nuclear data
library.

The present chapter provides the common ground: the transport equation from
which the multi-group description descends, the definition of the group
constants, the reason why that definition cannot be applied exactly, and the
error that follows from it.

\section{The neutron transport equation}

The stationary state of a multiplying system is described by the linear
Boltzmann equation for the angular flux $\phi(\vec{r},E,\vec{\Omega})$, the
density of neutrons at the position $\vec{r}$, with energy $E$, travelling in
the direction $\vec{\Omega}$, multiplied by their speed. In the form in which
criticality is expressed as an eigenvalue problem the equation reads
\cite{SPH_book_Bell_Glasstone}
%
\begin{align}
  \vec{\Omega}\cdot\nabla\phi(\vec{r},E,\vec{\Omega})
  &+ \Sigma_t(\vec{r},E)\,\phi(\vec{r},E,\vec{\Omega}) = \notag \\
  &\int_{0}^{\infty}\!\!\mathrm{d}E'\int_{4\pi}\!\!\mathrm{d}\vec{\Omega}'\,
   \Sigma_s(\vec{r},E'\!\to\!E,\vec{\Omega}'\!\to\!\vec{\Omega})\,
   \phi(\vec{r},E',\vec{\Omega}') \notag \\
  &+ \frac{\chi(E)}{4\pi\,k_{\mathrm{eff}}}
   \int_{0}^{\infty}\!\!\mathrm{d}E'\,
   \nu\Sigma_f(\vec{r},E')\,\Phi(\vec{r},E'),
  \label{mg:eq:transport}
\end{align}
%
where $\Sigma_t$ is the total macroscopic cross section, $\Sigma_s$ the
double-differential scattering cross section, $\nu\Sigma_f$ the production
cross section, $\chi$ the fission spectrum, and
%
\begin{equation}
  \Phi(\vec{r},E) = \int_{4\pi}\phi(\vec{r},E,\vec{\Omega})\,\mathrm{d}\vec{\Omega}
  \label{mg:eq:totalflux}
\end{equation}
%
the total flux, that is the angular flux integrated over the whole
solid angle. The two terms on the left describe the neutrons that leave a
point of the phase space, by streaming out of the volume element and by
interacting; the two on the right describe those that enter it, by scattering
from another energy and another direction and by fission. The multiplication
factor $k_{\mathrm{eff}}$ is the eigenvalue that makes the balance stationary,
and the fission source is written as isotropic, which the emission of fission
neutrons is to a very good approximation.

Equation~(\ref{mg:eq:transport}) has seven independent variables, three of
space, two of angle, one of energy and, in the time-dependent case, one of time.
A Monte Carlo code does not discretize any of them: it samples neutron histories
governed by the same cross sections and estimates the quantities of interest as
averages over the histories. A deterministic code discretizes all of them, and
the discretization of the energy variable is what this Part is concerned with.

\section{The multi-group form}\label{mg:sec:multigroup}

\subsection{Condensation in energy}

The energy interval of interest is partitioned into $G$ intervals, conventionally
numbered from the highest energy downwards, group $g$ covering
$E_{g}<E<E_{g-1}$. The group flux is the integral of the total flux over the
interval,
%
\begin{equation}
  \Phi_g(\vec{r}) = \int_{E_g}^{E_{g-1}}\Phi(\vec{r},E)\,\mathrm{d}E .
  \label{mg:eq:groupflux}
\end{equation}
%
Integrating Eq.~(\ref{mg:eq:transport}) over the same interval returns an
equation of identical structure, written in terms of $\Phi_g$ and of a set of
constants, provided that each constant is defined as the average of the
continuous cross section weighted by the flux,
%
\begin{equation}
  \Sigma_{x,g}(\vec{r}) =
  \frac{\displaystyle\int_{E_g}^{E_{g-1}}\Sigma_x(\vec{r},E)\,\Phi(\vec{r},E)\,\mathrm{d}E}
       {\displaystyle\int_{E_g}^{E_{g-1}}\Phi(\vec{r},E)\,\mathrm{d}E} ,
  \label{mg:eq:condensation}
\end{equation}
%
for a generic reaction $x$, and that the transfer between two groups is defined
by the corresponding double integral,
%
\begin{equation}
  \Sigma_{s,g'\to g}(\vec{r}) =
  \frac{\displaystyle\int_{E_g}^{E_{g-1}}\!\!\mathrm{d}E
        \int_{E_{g'}}^{E_{g'-1}}\!\!\mathrm{d}E'\,
        \Sigma_s(\vec{r},E'\!\to\!E)\,\Phi(\vec{r},E')}
       {\displaystyle\int_{E_{g'}}^{E_{g'-1}}\Phi(\vec{r},E')\,\mathrm{d}E'} .
  \label{mg:eq:scattering}
\end{equation}
%
The logic of Eq.~(\ref{mg:eq:condensation}) is contained in what happens when it
is multiplied by the group flux. The product $\Sigma_{x,g}\Phi_g$ is by
construction equal to $\int_g \Sigma_x(E)\Phi(E)\,\mathrm{d}E$, that is, to the
reaction rate of the continuous problem inside the interval. A group constant is
not an average cross section in any other sense: it is the number which, used
with the group flux, reproduces a reaction rate. The same reading applies to
Eq.~(\ref{mg:eq:scattering}) for the rate at which neutrons are transferred from
one group to another.

\subsection{Homogenization in space}

A deterministic solver of the kind used in this Part works on regions inside
which the composition is uniform, which the physical core is not: a fuel
assembly contains fuel, cladding, coolant and structural material, and the
solver is given one set of constants for the whole of it. The step that produces
them is the same average as before, performed over the volume $V_i$ of the
region instead of over the energy interval,
%
\begin{equation}
  \Sigma_{x,g,i} =
  \frac{\displaystyle\int_{V_i}\Sigma_{x,g}(\vec{r})\,\Phi_g(\vec{r})\,\mathrm{d}\vec{r}}
       {\displaystyle\int_{V_i}\Phi_g(\vec{r})\,\mathrm{d}\vec{r}} ,
  \label{mg:eq:homogenization}
\end{equation}
%
and it preserves, by the same argument, the reaction rate integrated over the
region. Condensation and homogenization are two instances of one operation, a
flux-weighted average of the cross section over a domain that the deterministic
model will not resolve, and in practice they are performed together in a single
Monte Carlo calculation which tallies the numerator and the denominator of
Eqs.~(\ref{mg:eq:condensation}) to (\ref{mg:eq:homogenization}) directly.

\subsection{The weighting flux}

The three definitions share a structure and, with it, a difficulty. The weight
is the flux, which is the unknown of the very problem the constants are meant to
solve. Were the true flux available, the multi-group equations would reproduce
the reaction rates of the continuous problem exactly, and there would be no
reason to solve them. What is available in its place is the flux of a model
which resembles the target: an infinite lattice of identical assemblies, an
assembly with a prescribed leakage, or, as in this Part, a full-core Monte Carlo
calculation of the system in one of its configurations.

The consequence is that the group constants are exact for the state in which
they were generated and approximate everywhere else. This is not a defect of the
procedure but its definition, and the practical question is not how to avoid the
approximation but how far the state of generation may differ from the state of
use before the error becomes unacceptable.

Two routes lead to that flux. The first is a dedicated deterministic code, which
solves the transport equation over a lattice cell or an assembly on a fine energy
mesh, with an explicit treatment of the resonance self-shielding: ECCO, the cell
code of the ERANOS system~\cite{eranos}, and MC$^2$-3~\cite{mc23} belong to the
fast-reactor tradition, DRAGON~\cite{dragon} to the thermal one. The second, the
one followed in this Part, is a Monte Carlo calculation of the system itself,
which tallies the numerator and the denominator of
Eqs.~(\ref{mg:eq:condensation}) to (\ref{mg:eq:homogenization}) with no
approximation on the geometry and none on the energy dependence, at the price of
a statistical uncertainty on every constant and of a cost that grows with the
number of states the library has to cover. The two routes are not alternatives in
principle, and which of them is preferable is decided by that number, as
Chapter~\ref{ch:sph-lfr} shows for the design studied here.

Whichever route produces it, the weighting flux need not be the one of a static
criticality calculation. That is the standard choice, and the one adopted in this
Part, but it remains a choice, and alternatives have been investigated.
\citet{Dugan2016,Dugan2018} weight with the eigenfunctions of the time-dependent
problem, so that the constants carry information on the free evolution of the
system, and report an improvement in the accuracy of transient calculations; the
same line has been followed at Politecnico di Torino, both on the spectral
formulations of the transport equation as weighting functions~\cite{abrateThesis}
and on the effect that the collapsing has on the effective neutron
lifetime~\cite{collapsing}.
% CV: qui vanno le due citazioni che mi mancano e che solo tu puoi darmi: il
% lavoro M&C 2023 fatto con Abrate, e il suo proseguimento a PHYSOR 2026 con
% Valeria e Martin che usa SCONE. Mandami autori, titolo esatto, conferenza,
% luogo e anno e le inserisco al posto di questa nota. La voce bib di SCONE
% (chiave "scone") e' gia' pronta in Chapter2/refs_mg.bib.
This Part does not act on the definition of the weighting
function: it takes the standard one and works on the two other degrees of
freedom, the correction of the constants that it produces and the placement of
the group boundaries over which the average is taken.

\section{Sources of error in the multi-group model}\label{mg:sec:errors}

A discrepancy between a deterministic result and a Monte Carlo reference
contains three contributions of different origin, and attributing it requires
that they be kept apart.

The first is the weighting spectrum and the environment in which it was
computed. What Eqs.~(\ref{mg:eq:condensation}) and (\ref{mg:eq:homogenization})
preserve are the reaction rates of the model used for the generation. When the
deterministic calculation places the same region in a different neighbourhood,
or moves a control rod, or reaches a different point of the cycle, the spectrum
inside the region changes and the constants are no longer the correct averages.
The error is concentrated where the spectrum varies most sharply in space, which
is at the interfaces between dissimilar regions and around the absorbers.
The second is the energy mesh. Inside a group the shape of the flux is lost, and
only its integral survives. A finer structure reduces the influence of the
assumed shape, since in the limit of an infinite number of groups
Eq.~(\ref{mg:eq:condensation}) becomes an identity, but the number of groups is
bounded by the cost of every calculation that will use them, and the benefit of
adding one is not uniform. What decides the accuracy is where the boundaries are
placed with respect to the features of the spectrum and of the cross sections,
and not the number of intervals alone \cite{collapsing}.

The third is the solver. Diffusion in place of transport, the spatial
discretization, the treatment of the anisotropy of scattering and of the
boundary conditions all contribute, and they do so whatever the quality of the
data. This term is not addressed in this Part. It is recalled here because the
comparisons reported in the following chapters contain it, and because a
correction applied to the data can, and sometimes does, compensate an error that
belongs to the solver.

Two features of a fast system give these contributions a particular weight.
The neutron spectrum extends over the whole energy range without the thermal
peak that anchors the structure of a light-water lattice, so that the placement
of the group boundaries has more freedom and less guidance. And the mean free
path is long compared with the pitch of an assembly, so that the flux inside a
region is determined by a wide neighbourhood, and the environment in which the
constants were generated matters more than in a thermal system.

\section{Equivalence and energy structure}\label{mg:sec:approaches}

Neither of the works that follow modifies Eq.~(\ref{mg:eq:transport}) or the
definitions of Sec.~\ref{mg:sec:multigroup}. Both act on the data that the
deterministic solver receives, and they act on the first two of the three
contributions above.

The first approach accepts that the constants cannot be exact and corrects them
so that the deterministic solution reproduces the reference reaction rates. This
is the idea of an equivalence procedure, of which superhomogenization (SPH)
\cite{Hbert_1980} is the instance adopted here: each homogenized region receives,
in each group, a multiplicative factor applied to its constants, and the factors
are computed iteratively until the deterministic model returns the reaction
rates of the Monte Carlo reference. The equivalence is imposed on a model of
limited extent and then transferred to the full core.
Chapter~\ref{ch:sph-lfr} develops the procedure for the Westinghouse LFR design,
compares two strategies for generating the constants, and examines what the
correction does to the multiplication factor, to the power distribution and to
the kinetic parameters.

The second approach acts on the energy mesh, treating the position of the group
boundaries as a decision variable rather than as a convention inherited from
practice. The structure is chosen from the internal boundaries of a fine mesh,
which makes the problem combinatorial, and the quantity to be optimized is the
agreement of a set of integral parameters with the Monte Carlo reference over
several configurations of the control system at once. Chapter~\ref{ch:energy-grid}
formulates that search, replaces the reference solver inside it with a neural
metamodel so that the campaign can be repeated many times, and recomputes every
structure it proposes with the solver, so that the conclusions rest on the
solver and not on the metamodel.

The two act on different terms and are complementary rather than alternative:
one repairs the constants of a region whose environment is not the one assumed,
the other chooses the intervals over which the averaging is performed. Applying
both to the same design would require a single calculation chain and is not
attempted here.
